% ---
% id: protocol-compiler/rem:relabel
% kind: rem
% title: "The setting-dependent relabelling after the chain"
% summary: >-
%   A choice whose bias depends on a station's setting is applied after the
%   chain, to both sides of the last link's relation, by one stochastic matrix;
%   by linearity the relation keeps the same channel. It is not a link, and it
%   is read through the setting alone.
% prominence: standard
% status: accepted
% uses: [def:weighted, def:compile, thm:choice-encoding, def:chaincond]
% mentions: [lem:weighted, rem:alphabet, reductions/lem:linear, reductions/rem:intertwining, reductions/def:joint-family, reductions/def:link, reductions/def:chain, reductions/ass:tsm, req:no-fixed-fraction, req:setting-only-weighting, req:choice-weighting-record]
% assumes: [A-COIN]
% refs: []
% provenance: newly-derived
% lean:
% ---

\begin{remark}\label{rem:relabel}
A public choice whose bias depends on the private setting of the apparatus
that measures a system, whichever party draws it --- a test or
generation round drawn against a probe setting, a probe light that fires
only in test rounds --- is the case the setting-dependent form of
Definition~\ref{def:weighted} exists for, and this remark records how the
compiler handles it. In the general format everything below holds per
measured factor, with that factor's receiver segment, family and setting map
in place of $\Chrec$, $\{F_v\}$ and $t$ (Definition~\ref{def:compile}).

\emph{What is applied, and where.} The compiler applies the stochastic matrix
$P(f\,|\,t)$, read through the setting map $t:\Vset\to\mathcal{T}$, to the
family $\{F_v\}$ the last link of $\Chrec$ delivered, \emph{after} the
reduction chain has been consumed and its
\tref[interface conditions on a chain]{interface conditions} decided
(Definition~\ref{def:chaincond}) and \emph{before} the validator and the
gates read the emitted declaration: the emitted family is
$N_{v,f}=P(f\,|\,t(v))F_v$ over the support $\Vset_P$, the emitted alphabet is
$\Vset_P$ in its product order, and the public part is declared on it
(Definition~\ref{def:compile}). The same weighting is, in the same stroke,
applied to the physical side of the last link's relation: the family the
emitted one is a reduction of is
$\{P(f\,|\,t(\gamma(w)))\Fph_w\}_{(w,f)}$ under the correspondence
$(w,f)\mapsto(\gamma(w),f)$ (Theorem~\ref{thm:choice-encoding}\,(vi)). Both sides are
relabelled by one stochastic matrix, which is what makes the next statement
true.

\emph{Why it is sound.} The soundness is
Theorem~\ref{thm:choice-encoding}\,(vi): by linearity of the adjoint, the
relation the last link asserts survives the stochastic matrix
$\delta_{vv'}P(f\,|\,t(v'))$ from $\Vset$ to $\Vset\times\mathcal{F}$, whose
columns sum to one, with the \emph{same} channel:
$\Qsys$, its sector map and its dimension are unchanged, every element
$P(f\,|\,t(v))F_v$ is a scalar multiple of an existing one so that block
diagonality holds by construction (Lemma~\ref{lem:weighted}\,(ii)), every
transported constraint operator has the same pull-back, and every entry's
certificate is untouched. That theorem is this module's instance, with the
one-line linearity argument written out, of the post-processing part~(i) of
\xref{reductions/lem:linear} and of the relabelling that
\xref{reductions/rem:intertwining} states and leaves to the consumer. What
that theorem does \emph{not}
assert is the choice marginal: the probability of $f$ is state-dependent
through the setting marginal of the outcome statistics
(Lemma~\ref{lem:weighted}\,(iv)), which is the output-state test-statistics
model (\xref{reductions/ass:tsm}, Requirement~\ref{req:no-fixed-fraction}),
and no fraction is factored out.

\emph{Why it is not a link.} A link in the sense of
\xref{reductions/def:link} composes a \emph{new} channel with the chain and
carries a \tref{reachable-set statement} for it; a relabelling composes no channel
--- it changes the alphabet at both ends of an existing relation and leaves
the channel where it was. It therefore has no kind, no register and no
position in a chain (\xref{reductions/def:chain}), and that module declares
none: the chain record this module consumes and re-emits is unchanged by the
weighting, the setting map it is read through is the class map of the joint
alphabet of \xref{reductions/def:joint-family}, and the weighting is recorded
beside the chain as this module's own construction
(Requirement~\ref{req:choice-weighting-record}). Putting $f$ into the alphabet
\emph{before} the chain instead is sound as well but strictly weaker where
the squasher is flag-shaped: every value of $f$ then opens flag sectors of its
own, and the squashed model gains adversarial freedom no physical state
produces; applied after the chain, no sector is indexed by $f$ and
$N_{v,f}=0$ as an operator wherever $P(f\,|\,t(v))=0$, so a zero-probability
pair is unreachable for every state rather than merely of zero weight. That
comparison is not a claim of this module; the claim is the soundness of the
relabelling, Theorem~\ref{thm:choice-encoding}\,(vi).

\emph{Why the weighting is read through the setting only.} $P$ is declared on
$\mathcal{T}\times\mathcal{F}$ and reaches an outcome through $t(v)$ alone.
A weighting that reached it through the device outcome would make $f$ an
inference from that outcome, not a choice, and the relabelling would then be
applied to a variable the adversary can influence through the state; that is
what Assumption~\ref{ass:coin} excludes and what
Requirement~\ref{req:setting-only-weighting} enforces by type. The setting
$t$ itself is a variable the squasher's channel depends on --- the conditional
families differ with it --- which is why it is in the alphabet before the
first link and $f$ is not.

\emph{With no setting-dependent choice declared}, the emitted family is the
chain's $\{F_v\}$ unchanged. What the emitted alphabet may be is the upstream
modules' question, not this module's (Remark~\ref{rem:alphabet}).
\end{remark}
